\problemname{Räkna öar}
\noindent

Visste du att Sverige är landet som har flest öar i världen? Det finns ca $267000$ stycken sammanlagt.
Många av dessa är dock väldigt små.
Kanske är det mer relevant att fråga vilket land som har flest öar om man bara räknar de öar med en landyta större än $100 m^2$, eller varför inte $1000 m^2$?

Det finns $N$ stycken länder, numrerade från $1$ till $N$.
Med hjälp av internet har du tagit reda på storleken av alla öarna i länderna.
För att tillfredställa din nyfikenhet vill du snabbt kunna hitta landet med flest öar om man bara räknar de som har minst $A$ kvadratmeter av landyta.
Du har $Q$ stycken värden på $A$ som du undrar över.

\section*{Indata}
Den första raden av indatan innehåller heltalen $N$ och $Q$ ($1 \leq N, Q \leq 10^5$),
antalet länder och antalet frågor.

De följande $N$ raderna beskriver vardera ett lands öar.
Varje av dessa rader börjar med ett heltal $c$ ($1 \leq c \leq 3 \cdot 10^5$),
antalet öar i det här landet.
Därefter följer $c$ heltal $s_1, \dots, s_c$ ($1 \leq s_i \leq 10^9$), storlekarna av landets öar. 

Till sist följer $Q$ rader som vardera innehåller ett heltal $A$ ($1 \leq A \leq 10^9$), vilket beskriver värdet av $A$ för en fråga.

Låt $T$ vara totala antalet öar, det vill säga är $T = \sum c$. Det håller att $T \leq 3 \cdot 10^5$.

\section*{Utdata}
För varje fråga, skriv ut ett heltal: numret på landet som har flest öar om man bara räknar de med minst $A$ kvadratmeter av landyta.
Om två länder har lika många öar så ska du svara med landet med högst nummer.

\section*{Poängsättning}
Din lösning kommer att testas på en mängd testfallsgrupper.
För att få poäng för en grupp så måste du klara alla testfall i gruppen.

\noindent
\begin{tabular}{| l | l | p{12cm} |}
  \hline
  \textbf{Grupp} & \textbf{Poäng} & \textbf{Gränser} \\ \hline
  $1$    & $10$       & $N, T, Q \leq 1000$ \\ \hline
  $2$    & $13$       & $s_i, A \leq 10$ \\ \hline
  $3$    & $16$       & $N \leq 2$ \\ \hline
  $4$    & $22$       & $N \leq 25$ \\ \hline
  $5$    & $39$       & Inga ytterligare begränsningar. \\ \hline
\end{tabular}

\section*{Förklaring av exempelfall $1$}
I första frågan i exempelfall $1$ har land $1$ tre öar med minst $1 m^2$ area: $3 7 9$.
Land $2$ har fyra öar med minst $1 m^2$ area: $1 2 2 3$.
Därmed har land $1$ flest öar.

I den andra frågan har land $1$ två öar som räknas, medan land $2$ inte har några öar alls.

\section*{Förklaring av exempelfall $2$}
I första frågan i exempelfall $2$ finns inga öar alls med area som minst $4 m^2$.
Därmed har alla länder samma antal öar och svaret är $2$, landet med höst nummer.
